\documentclass{beamer}

% Theme choice
\usetheme{Madrid}  % You can change the theme to other options like "Berlin", "Darmstadt", etc.

% Title page details
\title{Calcul différentiel }
\subtitle{Équation différentielle linéaire du premier ordre}
\author{SEDJRO }
\date{\today}
\input{commande_propre.sty}
\begin{document}





% Title page
\begin{frame}
    \titlepage
\end{frame}

% Table of contents
\begin{frame}
    \frametitle{Table des Matières}
    \tableofcontents
\end{frame}




\begin{frame}
\section{Introduction}
\frametitle{Introduction}
Souvent un problème concret fourni par l'analyse, la mécanique, la physique, la biologie , l'économie ,... se ramène à la résolution d'équations ou de systèmes différentielle. C'est dans cette otique que j'ai décide de vous faire un bref résumé de la résolution de ces gens de problème.
\end{frame}


%-----------------------------------------------------------------------------------------------

\newpage
\section{ Équation et système d'équations différentielle. }


% Section 1: Definitions and Examples
\subsection{Définitions et Exemples}

% Definition 2.1
\begin{frame}
    \frametitle{Définition 2.1}
    \begin{defi}
    On appelle équation linéaire du premier ordre une équation de la forme:
    \begin{equation*}\label{equ:1}
    \forall x \in \mathbb{R}, y'(x) + a(x)y(x) = f(x)
    \end{equation*}
    où $a$ et $f$ sont deux fonctions données, continues sur $I$, un intervalle non vide de $\mathbb{R}$. La fonction $a$ s'appelle le coefficient et la fonction $f$ le second membre de l'équation $(E)$.
    \end{defi}
\end{frame}

% Example 2.2
\begin{frame}
    \frametitle{Exemple 2.2}
    \begin{expl}
    L'équation suivante est une équation différentielle linéaire du premier ordre:
    \begin{equation*}
    \forall x \in \mathbb{R}, y'(x) + 2x y(x) = x.
    \end{equation*}
    \end{expl}
\end{frame}

% Definition 2.3
\begin{frame}
    \frametitle{Définition 2.3}
    \begin{defi}
    On appelle équation homogène associée à $(E)$ l'équation $(E_{H})$ obtenue en remplaçant $f$ par $0$ dans $(E)$:
    \begin{equation*}
    \forall x \in I, y'(x) + a(x)y(x) = 0.
    \end{equation*}
    \end{defi}
\end{frame}

% Proposition 2.4
\begin{frame}
    \frametitle{Proposition 2.4}
    \begin{prop}
    Si $y_{1}$ et $y_{2}$ sont deux solutions de l'équation $(E)$, alors $y_{1} - y_{2}$ est une solution de $(E_{H})$.
    \end{prop}
\end{frame}

% Theorem 2.5
\begin{frame}
    \frametitle{Théorème 2.5}
    \begin{theo}
    Soit $S_{H}$ l'ensemble des solutions de $E_{H}$, l'équation homogène associée à $(E)$, et soit $y_{p}$ une solution particulière de $(E)$. Alors l'ensemble $S$ des solutions de $(E)$ est défini par: 
    \[
    S = y_{p} + y_{h}, \quad \text{avec } y_{h} \in S_{H}.
    \]
    \end{theo}
\end{frame}

% Corollary 2.6
\begin{frame}
    \frametitle{Corollaire 2.6}
    \begin{coro}
    Soit $A$ une primitive quelconque de $a$ sur $I$ et soit $y_{p}$ une solution particulière de $(E)$. Alors, nous avons 
    \[
    S = \begin{cases}
    I \longrightarrow \mathbb{R} \\
    x \longmapsto y_{p}(x) + Ce^{-A(x)}
    \end{cases}
    \quad \text{avec } C \in \mathbb{R}.
    \]
    \end{coro}
\end{frame}

\end{document}
